\bookchapter{Κεφάλαιο 22: Πώς λειτουργεί η AI στο κινητό}

\chapterintro{Καλώς ήρθατε στην πιο ωραία “$\mu$αγική σκηνή” του κινητού, πρωταγωνιστής: η τεχνητή νοη$\mu$οσύνη. Ό$\mu$ορφες φωτογραφίες, φωνητικός βοηθός, πληκτρολόγιο που σας καταλαβαίνει — πίσω τους στέκει αυτός ο αόρατος εγκέφαλος.}

\booksection{I. AI Κά$\mu$ερα: το “darkroom” $\mu$ετά-επεξεργασίας $\mu$πήκε στο κινητό}

Η ση$\mu$ερινή AI φωτογραφία $\mu$οιάζει $\mu$ε ο$\mu$άδα επαγγελ$\mu$ατιών $\mu$έσα στο κλείστρο. Πριν πατήσετε, το AI δουλεύει: σαρώνει τη σκηνή, αναγνωρίζει “ουρανός ”, “πρόσωπο ”, “νύχτα”.

\booksection{II. Φωνητικός βοηθός: πώς σας καταλαβαίνει;}

\begin{enumerate}
  \item \textbf{Ακούει καθαρά (αναγνώριση φωνής)}: Κύ$\mu$ατα σε κεί$\mu$ενο.
  \item \textbf{Κατανοεί (Κατανόηση Φυσικής Γλώσσας)}. Κατανόηση πρόθεσης, εξαγωγή οντοτήτων.
  \item \textbf{Εκτέλεση και απάντηση}. Καλεί App/service, TTS $\mu$ιλά.
\end{enumerate}

\booksection{III. On-device AI vs. Cloud AI}

Το on-device AI κρατά τον υπολογισ$\mu$ό στο κινητό, γρήγορο, ιδιωτικό, offline. Το cloud AI θέλει δίκτυο αλλά χειρίζεται πιο περίπλοκες εργασίες.

\booksection{IV. Μεγάλα $\mu$οντέλα AI και κινητό: ο προσωπικός σας έξυπνος βοηθός}

2025-2026, AI(DeepSeek R1, ChatGPT-5,5.0,3). Έξυπνο.AI(NPU80 TOPS)τοπικό, AI.

\booksubsection{Τι είναι τα $\mu$εγάλα $\mu$οντέλα AI;}

Τα $\mu$εγάλα $\mu$οντέλα AI (LLM) είναι συστή$\mu$ατα AI εκπαιδευ$\mu$ένα $\mu$ε τεράστιους όγκους κει$\mu$ένων, ικανά να κατανοήσουν και να παράγουν φυσική γλώσσα. Μπορείτε να το φανταστείτε σαν έναν “σούπερ σπουδαστή” που έχει διαβάσει σχεδόν όλα τα κεί$\mu$ενα του διαδικτύου, ικανό να απαντά ερωτήσεις, να γράφει άρθρα, να $\mu$εταφράζει, να προγρα$\mu$$\mu$ατίζει και να αναλύει δεδο$\mu$ένα.

\booksubsection{Εφαρ$\mu$ογές $\mu$εγάλων $\mu$οντέλων AI στο κινητό:}

\begin{itemize}
  \item \textbf{Βοηθός συνο$\mu$ιλίας}: DeepSeek, ChatGPT, Doubao, Kimi: ερωτήσεις, $\mu$εταφράσεις, κεί$\mu$ενα, σύνοψη. Πιο έξυπνα από search — δίνουν απαντήσεις όχι links.
  \item \textbf{Γραφή και δη$\mu$ιουργία}. AI: emails, πρακτικά, βελτίωση, δη$\mu$ιουργία.
  \item \textbf{Δη$\mu$ιουργία και επεξεργασία εικόνων}: AI αφαιρεί περαστικούς, επεκτείνει άκρα, αλλάζει στυλ, δη$\mu$ιουργεί από κεί$\mu$ενο.
  \item \textbf{Ζωντανή $\mu$ετάφραση}. Live call, face-to-face, photo translate — ταξίδια χωρίς γλώσσα.
  \item \textbf{Εκπαιδευτική υποστήριξη}: AI εξηγεί έννοιες, δίνει ασκήσεις, διορθώνει — φορητός δάσκαλος.
\end{itemize}

\booksubsection{On-device $\mu$εγάλο $\mu$οντέλο: το κινητό τρέχει AI}

Τα νεότερα flagships τρέχουν τοπικά $\mu$ικρά AI $\mu$οντέλα (συνήθως \textasciitilde{}7B παρα$\mu$έτρους), φέρνοντας:

\begin{itemize}
  \item \textbf{Χωρίς δίκτυο}. Offline AI, αεροπλάνο, $\mu$ετρό.
  \item \textbf{πιο ασφαλής ιδιωτικότητα}: Οι συζητήσεις δεν ανεβαίνουν στο cloud, τοπική επεξεργασία.
  \item \textbf{Ταχύτερη απόκριση}. Χωρίς latency, AI ά$\mu$εση απάντηση.
  \item \textbf{Τρέχοντες περιορισ$\mu$οί}. On-device υπολείπεται cloud, περίπλοκα θέλουν δίκτυο. Αλλά chips βελτιώνονται.
\end{itemize}

\booksubsection{Ενσω$\mu$άτωση AI σε επίπεδο συστή$\mu$ατος:}

\begin{itemize}
  \item \textbf{Apple Apple Intelligence}: AI iOS 18, εργαλεία γραφής, magic wand, Siri, σύνοψη ειδοποιήσεων. On-device πρώτα, cloud για περίπλοκα.
  \item \textbf{Huawei Pangu LLM}. Σε HarmonyOS, Xiaoyi $\mu$ε LLM: summary, smart screen, AI editing.
  \item \textbf{Xiaomi HyperAI}. System AI: AI search, summary, live captions, on-device chat.
  \item \textbf{Samsung Galaxy AI}. Live call translate, AI notes, image gen/edit, circle search.
\end{itemize}

\booksection{V. Τυπικά σενάρια και πρακτική}

Για να νιώσετε τις “υπερδυνά$\mu$εις ” του AI, παρακάτω πέντε σενάρια δείχνουν πώς το AI “σώζει” την καθη$\mu$ερινότητα:

\booksubsection{Σενάριο 1: Ζωντανή $\mu$ετάφραση — ο προσωπικός διερ$\mu$ηνέας για ταξίδια στο εξωτερικό}

\textbf{Πραγ$\mu$ατική απόδοση}: Στο Παρίσι χωρίς γαλλικά; iPhone “Translate” face-to-face, subtitles, ακό$\mu$α και simultaneous. Samsung Galaxy AI: κλήση σε κινέζικα, ακούνε αγγλικά.

\textbf{Αρχή πίσω}. ASR NMT TTS, latency <1 δευτ.

\booksubsection{Σενάριο 2: AI αφαίρεση — σώστε τις τουριστικές φωτογραφίες που “καταστράφηκαν” από περαστικούς}

\textbf{Πραγ$\mu$ατική απόδοση}: Μια ωραία selfie $\mu$ε τουρίστες; “Magic eraser” κύκλος, διαγράφονται, συ$\mu$πληρώνεται, εκπληκτικό.

\textbf{Αρχή πίσω}: Inpainting: αναγνωρίζει και διαγράφει, $\mu$ετά “$\mu$αντεύει” συ$\mu$πλήρωση. Εκατο$\mu$$\mu$ύρια εικόνες εκπαίδευσης.

\booksubsection{Σενάριο 3: Επεξεργασία $\mu$ε AI — γίνετε “φωτογράφος” $\mu$ε ένα κλικ}

\textbf{Πραγ$\mu$ατική απόδοση}. Α$\mu$υδρή φωτογραφία; AI “magic sky” = clear, sunset, stars. AI outpainting από portrait σε landscape.

\textbf{Αρχή πίσω}: Generative AI Diffusion Model, κατανοεί και δη$\mu$ιουργεί pixels. Αιχ$\mu$ή.

\booksubsection{Σενάριο 4: Επεξεργασία εγγράφων — πολλαπλασιαστής απόδοσης για εργαζό$\mu$ενους γραφείου}

\textbf{Πραγ$\mu$ατική απόδοση}: Μετά τη σύσκεψη, δώστε ηχόγρα$\mu$$\mu$α σε AI: δο$\mu$η$\mu$ένα πρακτικά, tasks, σύνοψη. PDF 10 σελίδων σε 10 δευτ. βγάζει συ$\mu$περάσ$\mu$ατα.

\textbf{Αρχή πίσω}. Long Context LLM διαβάζει δεκάδες χιλιάδες λέξεις, multi-round reasoning. Apple Intelligence, Huawei Xiaoyi, WPS AI, Tencent Yuanbao.

\booksubsection{Σενάριο 5: Αναζήτηση εικόνων και αναγνώριση αντικει$\mu$ένων $\mu$ε AI — η “οπτική $\mu$ηχανή αναζήτησής” σας}

\textbf{Πραγ$\mu$ατική απόδοση}: Δείτε λουλούδι, AI λέει “hibiscus”. Στο άλ$\mu$που$\mu$ ψάξτε “θάλασσα”, AI βρίσκει.

\textbf{Αρχή πίσω}. CV models embedding, matching $\mu$ε κεί$\mu$ενο. “Multi-modal search” = άλ$\mu$α.

\booksection{VI. Βελτιστοποίηση prompt: ο “τύπος επικοινωνίας”}

Η ποιότητα των απαντήσεων του AI εξαρτάται από το πώς “ρωτάτε”. Το ίδιο ερώτη$\mu$α, $\mu$ε διαφορετική διατύπωση, δίνει εντελώς διαφορετικές απαντήσεις. Παρακάτω παρουσιάζου$\mu$ε ένα\textbf{η πιο πρακτική “ρόλος-ιστορικό-εργασία-περιορισ$\mu$οί” προσέγγιση}, $\mu$ε πολλά παραδείγ$\mu$ατα.

\begin{tcbinfo}{Γιατί χρειάζεται πλαίσιο;}
Πολλοί χρησι$\mu$οποιούν την AI όπως το Baidu — πληκτρολογούν $\mu$ερικές λέξεις-κλειδιά. Αλλά η AI δεν είναι $\mu$ηχανή αναζήτησης, είναι ένα γλωσσικό $\mu$οντέλο που “διαβάζει, κατανοεί και απαντά”.\textbf{δώστε περισσότερο πλαίσιο, ξεκαθαρίστε τι θέλετε, περιορίστε τις αποκλίσεις}, σας δίνει ακριβή αποτελέσ$\mu$ατα. Το πλαίσιο $\mu$εταφράζει “τι είναι καλή απάντηση” σε εντολές AI.

\end{tcbinfo}

\booksubsection{Ρόλος (Role): Πείτε στο AI ποιο είναι}

Δώστε στο AI συγκεκρι$\mu$ένη “ταυτότητα” και θα αντλήσει γνώσεις και ύφος του αντίστοιχου το$\mu$έα.\textbf{Όσο πιο συγκεκρι$\mu$ένος ο ρόλος, τόσο πιο επαγγελ$\mu$ατική η έξοδος}.

\booksubsection{Πρακτικό παράδειγ$\mu$α:}

\begin{itemize}
  \item "Είστε ένας\textbf{$\mu$ηχανικός επισκευής Xiaomi $\mu$ε 10 χρόνια ε$\mu$πειρίας}." (τεχνικό θέ$\mu$α)
  \item "Είστε ένας\textbf{καθηγητής φυσικής γυ$\mu$νασίου που απλοποιεί περίπλοκες έννοιες}." (επιστη$\mu$ονικό ερώτη$\mu$α)
  \item "Είστε ένας\textbf{Εκδότης βιβλίων $\mu$ε καλό γούστο, πολυδιαβασ$\mu$ένος, $\mu$ε δικό του στυλ}." (θέ$\mu$α γραφής)
  \item "Είστε ένας\textbf{food blogger που φτιάχνει πρωινό για εργαζό$\mu$ενους}." (θέ$\mu$α καθη$\mu$ερινότητας)
  \item "Είστε ένας\textbf{TERN, υπο$\mu$ονετικός γιατρός TCM που ποτέ δεν ειρωνεύεται αρχαρίους}." (θέ$\mu$α υγείας)
\end{itemize}

\textbf{Γιατί είναι αποτελεσ$\mu$ατικό;}Ο ρόλος ορίζει στο AI\textbf{εύρος γνώσης},\textbf{ύφος έκφρασης}και\textbf{επαγγελ$\mu$ατικό βάθος}. Ένα AI χωρίς ρόλο είναι “γνώστης όλων αλλά ειδικός σε τίποτα ”· $\mu$ε ρόλο, “εισέρχεται” στην ταυτότητα και σκέφτεται σαν επαγγελ$\mu$ατίας.

\booksubsection{Ιστορικό (Background): Πείτε στο AI την κατάστασή σας}

Δηλώστε ξεκάθαρα “ποιος είστε, ποια είναι η τρέχουσα κατάστασή σας, ποιοι είναι οι περιορισ$\mu$οί σας” — το AI $\mu$πορεί να δώσει\textbf{πραγ$\mu$ατικά κατάλληλο}απαντήσεις, όχι αόριστες.

\booksubsection{Πρακτικό παράδειγ$\mu$α:}

\begin{itemize}
  \item "το κινητό $\mu$ου είναι\textbf{Xiaomi 14, αγορασ$\mu$ένο πριν έξι $\mu$ήνες, ζεσταίνεται πολύ στα παιχνίδια}." (κατάσταση συσκευής)
  \item "εί$\mu$αι\textbf{$\mu$ια 60 χρονη συνταξιούχος, $\mu$ε όχι τόσο καλά $\mu$άτια, που δεν χρησι$\mu$οποιεί εύκολα το smartphone}." (προφίλ χρήστη)
  \item "αυτή την περίοδο\textbf{προετοι$\mu$άζο$\mu$αι για CET-6 Αγγλικών τον Ιούνιο 2026, $\mu$έτριο επίπεδο}." (εκπαιδευτικό ιστορικό)
  \item "άνοιξα ένα\textbf{$\mu$ικρό καφέ σε τριτεύουσα πόλη, $\mu$ηνιαίος τζίρος 30.000, δύο υπάλληλοι}." (επιχειρη$\mu$ατικό ιστορικό)
  \item "εγώ\textbf{Μόλις γέννησα και εί$\mu$αι στο σπίτι σε περίοδο λοχείας, $\mu$ε σοβαρή έλλειψη ύπνου}." (στάδιο ζωής)
\end{itemize}

\textbf{Γιατί είναι αποτελεσ$\mu$ατικό;}Το ιστορικό λέει στο AI\textbf{“για ποιον υπηρετεί”}και\textbf{“υπό ποιες συνθήκες”}. Το ίδιο ερώτη$\mu$α εξηγείται διαφορετικά σε $\mu$αθητή και γονιό· διαφορετικές συ$\mu$βουλές σε κατοίκους πόλης και αγροτικών περιοχών. Χωρίς ιστορικό, το AI δίνει απαντήσεις “ελάχιστου κοινού διαιρέτη”, συχνά αόριστες.

\booksubsection{Εργασία (Task): Πείτε ξεκάθαρα στο AI τι να κάνει}

Αυτό το βή$\mu$α θέλει\textbf{Ξεκινήστε $\mu$ε ρή$\mu$α}(γράψτε, καταγράψτε, συγκρίνετε, $\mu$εταφράστε, εξηγήστε, συνοψίστε...), περιγράψτε την εργασία\textbf{συγκεκρι$\mu$ένα, εκτελέσι$\mu$α}. Ασαφής εργασία φέρνει ασαφείς απαντήσεις.

\booksubsection{Πρακτικό παράδειγ$\mu$α:}

\begin{itemize}
  \item "παρακαλώ\textbf{Καταγράψτε 5 δι και ώ$\mu$ατα που πρέπει να κλείσω ά$\mu$εσα}, $\mu$ε αιτία. (ξεκαθαρισ$\mu$ένη ενέργεια)
  \item "παρακαλώ\textbf{$\mu$εταφράστε αυτό το αγγλικό email στα κινεζικά}, $\mu$ε επαγγελ$\mu$ατικό ύφος. (συγκεκρι$\mu$ένη εργασία)
  \item "παρακαλώ\textbf{Συγκρίνετε αυτά τα τέσσερα κινητά στην τι$\mu$ή των 3000 γιουάν}, $\mu$ε πίνακα πλεονεκτη$\mu$άτων-$\mu$ειονεκτη$\mu$άτων. ($\mu$ορφή)
  \item "παρακαλώ\textbf{εξηγήστε "τι είναι 5G" σε $\mu$αθητή δη$\mu$οτικού}." (ξεκαθαρισ$\mu$ένο στυλ)
  \item "παρακαλώ\textbf{γράψτε τρία σενάρια $\mu$ικρού βίντεο για το πρωινό $\mu$ου}, για εργαζό$\mu$ενους, ζευγάρια, γονείς. (ξεκαθαρισ$\mu$ένη κατάσταση)
\end{itemize}

\textbf{Γιατί είναι αποτελεσ$\mu$ατικό;}Η ενότητα εργασίας είναι το “κέντρο εντολών” του AI.\textbf{Όσο πιο σαφές το ρή$\mu$α, τόσο πιο συγκεκρι$\mu$ένη η απαίτηση, τόσο καλύτερα η AI γνωρίζει «τι να παραδώσει»}. Ασαφείς φράσεις όπως “καν' ένα ”, “φτιάξε ”, “δες” είναι τα $\mu$εγαλύτερα λάθη $\mu$ε το AI — θα $\mu$αντέψει και συχνά θα αποκλίνει.

\booksubsection{Περιορισ$\mu$οί (Constraints): Πείτε στο AI τι να $\mu$ην κάνει}

Το “όριο ” ακούγεται σαν “περιορισ$\mu$ός”, στην πραγ$\mu$ατικότητα είναι\textbf{να συγκλίνει το AI στην πιο κατάλληλη έξοδο}εργαλείο. Συνηθισ$\mu$ένες διαστάσεις: λέξεις, $\mu$ορφή, ύφος, αποφυγές, απαιτού$\mu$ενα.

\booksubsection{Πρακτικό παράδειγ$\mu$α:}

\begin{itemize}
  \item "\textbf{Όχι πάνω από 200 χαρακτήρες, σε ση$\mu$εία}." ($\mu$ήκος+$\mu$ορφή)
  \item "\textbf{Ύφος φιλικό, σαν συζήτηση φίλων}, χωρίς "πρώτον, δεύτερον" ύφος. (ύφος)
  \item "\textbf{Μην χρησι$\mu$οποιείτε τεχνικούς όρους}, εξηγώντας όρους. (περιορισ$\mu$ός λέξεων)
  \item "\textbf{Κάθε παράγραφος έως 3 γρα$\mu$$\mu$ές}, για ανάγνωση σε κινητό. (διάταξη)
  \item "\textbf{Πρέπει να περιλα$\mu$βάνει: σενάρια χρήσης, πλεονεκτή$\mu$ατα/$\mu$ειονεκτή$\mu$ατα, αντιστοιχία $\mu$ε τον προϋπολογισ$\mu$ό $\mu$ου}." (απαιτού$\mu$ενα στοιχεία)
  \item "\textbf{Μην δίνετε συγκεκρι$\mu$ένες επενδυτικές συ$\mu$βουλές}, $\mu$όνο εξήγηση αρχών. (όριο)
\end{itemize}

\textbf{Γιατί είναι αποτελεσ$\mu$ατικό;}Χωρίς περιορισ$\mu$ούς, το AI “αποκλίνει” — πολύ ακαδη$\mu$αϊκό, πολύ $\mu$εγάλο ή εκτός θέ$\mu$ατος. Σαφείς περιορισ$\mu$οί = σαφές “πρότυπο παράδοσης”.

\booksubsection{Αξιολόγηση ποιότητας prompt: 5 χαρακτηριστικά καλού prompt}

Το να γράφετε καλά prompts δεν είναι εφάπαξ,\textbf{$\mu$ετά γράψι$\mu$ο ελέγξτε αυτά τα 5}, πριν στείλετε στο AI:

\begin{enumerate}
  \item \textbf{συγκεκρι$\mu$ένα, όχι αφηρη$\mu$ένα}: όχι "γράψε άρθρο ", αλλά "800 λέξεις, 3 ενότητες, για νέες $\mu$α$\mu$άδες, οδηγός αγοράς".
  \item \textbf{Περιλα$\mu$βάνει φόντο, όχι απο$\mu$ονω$\mu$ένα}: Πείτε στο AI "ποιος εί$\mu$αι, γιατί ρωτάω ", δίνει "customized" απαντήσεις.
  \item \textbf{ξεκαθαρισ$\mu$ένη $\mu$ορφή εξόδου}. Table, list, paragraph, JSON; πείτε πριν.
  \item \textbf{επαρκείς περιορισ$\mu$οί}: Λέξεις, ύφος, αποφυγές, ένα από τα τρία = βελτίωση.
  \item \textbf{Μπορεί να επαληθευτεί από AI}. Ζητήστε από το AI να επαναλάβει· αν δεν $\mu$πορεί, δεν ήταν ξεκάθαρο.
\end{enumerate}

\booksubsection{Συχνά προβλή$\mu$ατα και λύσεις}

\begin{center}
\begin{tabularx}{\linewidth}{|X|X|}
\hline
\textbf{Συ$\mu$πτώ$\mu$ατα} & \textbf{Λύση} \tabularnewline
\hline
Η απάντηση είναι πολύ γενική, $\mu$η πρακτική & Προσθέστε πλαίσιο (ποιος είστε / σε ποιο σενάριο) \tabularnewline
Πολύ ακαδη$\mu$αϊκή/υπερβολικά επίση$\mu$η απάντηση & Προσθέστε περιορισ$\mu$ό τόνου + λίστα απόρριψης «επίση$\mu$ης γλώσσας» \tabularnewline
Η απάντηση ξεφεύγει από το θέ$\mu$α/δεν απαντά στην ερώτηση & Επαναλάβετε την εργασία $\mu$ε ρή$\mu$α και ζητήστε «απαντήστε $\mu$όνο σε αυτό» \tabularnewline
Η απάντηση είναι πολύ $\mu$εγάλη για να διαβαστεί & Αυστηρό όριο λέξεων ($\mu$έσα σε 200) \tabularnewline
Το AI $\mu$ιλάει σοβαρά αλλά λέει ανοησίες & Προσθέστε "Εάν δεν είστε σίγουροι, ση$\mu$ειώστε 'δεν εί$\mu$αι σίγουρος'" + αυτοέλεγχος ση$\mu$αντικών γεγονότων \tabularnewline
Δεν αρέσει η πρώτη έκδοση & Μην ξαναγράψετε, προσθέστε "λίγο πιο..." στο αρχικό prompt \tabularnewline
\hline
\end{tabularx}
\end{center}

\booksubsection{Τρία πλήρη παραδείγ$\mu$ατα σεναρίων}

\begin{tcbinfo}{Σενάριο 1: Βοηθός καθη$\mu$ερινότητας (προτάσεις πρωινού)}
\textbf{Prompt (προτροπή)}:

Είσαι food blogger ειδικός στη διατροφή (\textbf{Ρόλος}). Εί$\mu$αι εργαζό$\mu$ενη $\mu$ητέρα $\mu$ε $\mu$όνο 15 λεπτά κάθε πρωί για πρωινό, $\mu$ε 6 χρονο παιδί και πεθερά 60 ετών (\textbf{Πλαίσιο}). Παρακαλώ\textbf{καταγράψτε 7 $\mu$έρες χωρίς επανάληψη}(\textbf{Εργασία}). Απαιτήσεις: (1) $\leq$15 λεπτά (2) Ισορροπία (3) Όλες οι ηλικίες (4) Σούπερ $\mu$άρκετ (\textbf{Περιορισ$\mu$οί}).

\textbf{Γιατί λειτουργεί}: Ρόλος=ειδίκευση, πλαίσιο=3 τύποι, εργασία=παραδοτέα, 4 περιορισ$\mu$οί.

\end{tcbinfo}

\begin{tcbinfo}{Σενάριο 2: Εκπαίδευση (CET-6)}
\textbf{Prompt (προτροπή)}:

Είσαι\textbf{έ$\mu$πειρος καθηγητής CET-6}(ρόλος). Προπτυχιακός,\textbf{CET-4 480, στόχος CET-6 500, 3 $\mu$ήνες}(πλαίσιο). Παρακαλώ\textbf{δη$\mu$ιουργήστε 12-εβδο$\mu$άδων σχέδιο}(εργασία). (1) Χρόνος/εβδο$\mu$άδα (2) Κατεύθυνση/$\mu$ήνα (3) 3 βιβλία (4) Πίνακας.

\textbf{Γιατί λειτουργεί}: Ρόλος=υπόβαθρο, πλαίσιο=αναφορές, εργασία=“φάσεις ”, περιορισ$\mu$ός=“πίνακας”.

\end{tcbinfo}

\begin{tcbinfo}{Σενάριο 3: Αποδοτικότητα (πρακτικά)}
\textbf{Prompt (προτροπή)}:

Είσαι\textbf{έ$\mu$πειρος διοικητικός βοηθός}(ρόλος). Σύναψη 1 ώρας,\textbf{συζήτηση $\mu$άρκετινγκ Q2}(πλαίσιο).\textbf{οργανώστε πρακτικά}(εργασία). (1) 3 $\mu$έρη (2) Υπεύθυνος+προθεσ$\mu$ία (3) $\leq$600 (4) markdown.

\textbf{Γιατί λειτουργεί}: Ρόλος=διοικητική οπτική, πλαίσιο=θέ$\mu$α, εργασία=δο$\mu$ή, 4 περιορισ$\mu$οί.

\end{tcbinfo}

\textbf{Προχωρη$\mu$ένα: AI ως “βοηθός”}\par

Δοκι$\mu$άστε προχωρη$\mu$ένες χρήσεις:

\begin{itemize}
  \item \textbf{Πολλαπλοί γύροι}: "Σαν σε δη$\mu$οτικό".
  \item \textbf{Chain of Thought (CoT)}: "Βή$\mu$α-βή$\mu$α σκέψη".
  \item \textbf{Πολλαπλά AI}: Α γράφει, Β επι$\mu$ελείται, Γ $\mu$εταφράζει.
  \item \textbf{System prompt}: Προετοι$\mu$άστε ταυτότητα/προτι$\mu$ήσεις.
\end{itemize}

\booksection{Όρια AI: Διατηρήστε κρίση}

Ισχυρό αλλά όχι παντοδύνα$\mu$ο. Αναγνωρίστε όρια:

\begin{itemize}
  \item \textbf{"Ψευδαισθήσεις"}: Το AI επινοεί γεγονότα. Επαληθεύστε.
  \item \textbf{Cutoff γνώσης}: Τα περισσότερα έχουν cutoff.
  \item \textbf{Περιορισ$\mu$ένος συλλογισ$\mu$ός}: Κάνει λάθη σε πολύπλοκα.
  \item \textbf{Κίνδυνος απορρήτου}: Cloud AI = κίνδυνος. Ευαίσθητα = on-device.
  \item \textbf{Προκαταλήψεις}: Φέρει προκαταλήψεις.
\end{itemize}

\begin{tcbinfo}{Συ$\mu$βουλές AI}
\begin{itemize}
  \item Μπορεί να είναι ανακριβές.
  \item Μην δίνετε προσωπικά δεδο$\mu$ένα.
  \item Καλό prompt = υψηλότερη ποιότητα.
  \item Εργαλείο όχι αυθεντία.
  \item Σε κρίσι$\mu$ες αποφάσεις, ειδικοί.
\end{itemize}

\end{tcbinfo}

\begin{tcbpractice}{Πρόκληση: 3 AI ασκήσεις}
\begin{enumerate}
  \item \textbf{Πρόκληση 1: Σύνθετη εντολή}: Πείτε εντολή πολλαπλών βη$\mu$άτων.
  \item \textbf{Πρόκληση 2: AI editing}: Αφαιρέστε περαστικούς.
  \item \textbf{Πρόκληση 3: Prompt}: Απλό vs CRISPE.
\end{enumerate}

\end{tcbpractice}

\begin{tcbinfo}{Περίπτωση 2: Πλήρης ροή $\mu$εταφοράς δεδο$\mu$ένων αλλαγής κινητού}
\begin{itemize}
  \item \textbf{Σενάριο}: iPhone 13 Xiaomi 15 αλλαγή $\mu$άρκας.
  \item \textbf{Βή$\mu$ατα}:(1) iCloud/αντίγραφο ασφαλείας (2) Apple ID + απενεργοποίηση "εύρεση" (3) ""App (4) $\mu$εταφορά (5) "καταγραφή $\mu$εταφορά" (6) επαλήθευση Δεδο$\mu$ένα(//) (7) Κρυπτογράφηση.
  \item \textbf{Ανα$\mu$ενό$\mu$ενο αποτέλεσ$\mu$α}: Μηδενική απώλεια δεδο$\mu$ένων, παλιό κινητό ασφαλής $\mu$εταπώληση.
\end{itemize}

\end{tcbinfo}
